\documentclass[10pt]{article}
\usepackage[T1]{fontenc}
\usepackage[utf8]{inputenc}
\usepackage[french]{babel}
\usepackage{newpxtext,newpxmath}
\usepackage[a4paper,top=1.5cm,bottom=1.35cm,left=1.6cm,right=1.6cm]{geometry}
\usepackage{microtype}
\usepackage{xcolor}
\usepackage{titlesec}
\usepackage[most]{tcolorbox}
\usepackage{enumitem}
\usepackage{array}
\usepackage{booktabs}
\usepackage{fancyhdr}

\definecolor{navy}{HTML}{1F3A5F}
\definecolor{grey}{HTML}{54606E}
\definecolor{green}{HTML}{2C6E49}
\definecolor{amber}{HTML}{9C6A16}
\definecolor{red}{HTML}{A6322E}
\definecolor{plum}{HTML}{7A4EA8}
\definecolor{boxgray}{HTML}{F3F5F8}
\definecolor{boxred}{HTML}{FCF2F1}
\definecolor{boxgreen}{HTML}{F1F7F3}
\definecolor{boxplum}{HTML}{F7F3FB}
\definecolor{boxcode}{HTML}{F7F8FA}
\definecolor{boxedge}{HTML}{D9E0EA}

\setlength{\parindent}{0pt}
\setlength{\parskip}{2.6pt}
\linespread{1.0}

\pagestyle{fancy}\fancyhf{}
\lhead{\small\color{grey}Scella · architecture}
\rhead{\small\color{grey}MZ --- La mesure du Cadrage --- 23 août 2026}
\cfoot{\small\color{grey}\thepage}
\renewcommand{\headrulewidth}{0.3pt}

\titleformat{\section}{\color{navy}\large\bfseries}{}{0pt}{}
\titlespacing{\section}{0pt}{11pt}{3pt}
\titleformat{\subsection}{\color{navy}\bfseries}{}{0pt}{}
\titlespacing{\subsection}{0pt}{7pt}{1.5pt}

\newtcolorbox{these}{colback=boxgray,colframe=boxedge,boxrule=0.5pt,arc=2.5pt,
  left=8pt,right=8pt,top=4.5pt,bottom=4.5pt}
\newtcolorbox{alerte}{colback=boxred,colframe=red!55,boxrule=0.8pt,arc=2.5pt,
  left=8pt,right=8pt,top=4.5pt,bottom=4.5pt}
\newtcolorbox{acquis}{colback=boxgreen,colframe=green!45,boxrule=0.8pt,arc=2.5pt,
  left=8pt,right=8pt,top=4.5pt,bottom=4.5pt}
\newtcolorbox{piste}{colback=boxplum,colframe=plum!45,boxrule=0.8pt,arc=2.5pt,
  left=8pt,right=8pt,top=4.5pt,bottom=4.5pt}
\newtcolorbox{extrait}{colback=boxcode,colframe=boxedge,boxrule=0.5pt,arc=1.5pt,
  left=7pt,right=7pt,top=3.5pt,bottom=3.5pt}

\newcommand{\best}[1]{\textcolor{green}{\textbf{#1}}}
\newcommand{\pire}[1]{\textcolor{red}{\textbf{#1}}}
\newcommand{\sz}[1]{\textcolor{grey}{\small\textsf{#1}}}
\newcommand{\code}[1]{\texttt{\small #1}}
\newcommand{\fl}{$\rightarrow$}
\newcommand{\mm}[1]{\textbf{\textcolor{navy}{#1}}}

\begin{document}

\begin{center}
{\Huge\bfseries\color{navy} Scella}\\[4pt]
{\large\color{grey}\itshape MZ --- La mesure du Cadrage · architecture}\\[4pt]
{\small\color{grey} 23 août 2026 \; · \; hors chemin \; · \; indépendant, local, \textbf{gratuit}}
\end{center}
\vspace{2pt}

\begin{these}
\small \mm{MZ} n'est pas une fonctionnalité : c'est une \textbf{mesure}, et la feuille de route dit
qu'elle « peut tuer le Cadrage sans rien coûter ». Un module dont le produit attendu est peut-être un
\textbf{refus} se conçoit autrement qu'un module qui livre : ses critères de réfutation s'écrivent
\emph{avant} l'expérience, et sa réussite se juge à sa capacité de dire non.

Ce document établit d'abord ce que la lecture du code apprend --- trois faits qui changent la
question posée --- puis définit les quantités, le protocole, et \textbf{ce qui tue le Cadrage},
déclaré d'avance.
\end{these}

\section{1 · Ce que le Cadrage promet, et ce que le code en dit}

\subsection{1.1 · La promesse}

Le Cadrage --- le « pré-mâchage » --- fait passer l'entrée par un petit modèle local
(\code{compressor.gguf}, Llama-3.2-1B Q5, 911\,Mo) avant de l'envoyer au modèle frontière. Le type
qui le décrit annonce sa thèse en une ligne :

\begin{extrait}\small
\code{//! Compresse un texte long \textbf{en préservant des entités spécifiées} et en}\\
\code{//! produisant un ratio mesurable (compression typique \textbf{8:1 à 12:1}).}
\end{extrait}

Deux promesses, donc, et non une : un \textbf{ratio} et une \textbf{préservation}. L'économie vient
de la première ; la sûreté d'employer le procédé vient de la seconde.

\subsection{1.2 · Ce qui a déjà été mesuré, et ce que cela ne dit pas}

Un rapport existe (\code{CHEWING\_SAVINGS\_REPORT\_v2.md}, 30 mai 2026) :

\begin{center}\small
\renewcommand{\arraystretch}{1.2}
\begin{tabular}{@{}>{\raggedright\arraybackslash}p{5.2cm} >{\raggedright\arraybackslash}p{3.2cm} >{\raggedright\arraybackslash}p{3.2cm} >{\raggedright\arraybackslash}p{2.6cm}@{}}
\toprule
& \textbf{Baseline} & \textbf{Full (S4+S6)} & \textbf{Écart} \\
\midrule
Coût total (6 fixtures) & 2{,}426 ¢ & 2{,}246 ¢ & \best{-7 \%} \\
Latence moyenne & 5\,282 ms & 9\,473 ms & \pire{+79 \%} \\
Modes de panne observés & 0 & 4 & \pire{+4} \\
\bottomrule
\end{tabular}
\end{center}

Ce rapport mesure un \textbf{pipeline de bout en bout}, sur six fixtures, avec de vraies API. Il ne
dit rien du \textbf{levier} lui-même : ni jusqu'où le compresseur peut aller, ni ce qu'il perd en
chemin. C'est exactement la distinction que \mm{HB.3} et \mm{HB.4} ont payée ailleurs dans ce
dépôt --- mesurer le levier \emph{avant} de bâtir dessus, faute de quoi on optimise un pipeline
autour d'un composant dont personne ne connaît le plafond.

\subsection{1.3 · Trois faits établis en lisant le code}

\begin{alerte}
\textbf{F1 --- La préservation n'est jamais demandée.} Les \textbf{deux seuls} sites de production
qui appellent le compresseur (\code{mission\_a.rs:398}, \code{mission\_b.rs:219}) passent
\code{must\_preserve: vec![]}. Le $-7\,\%$ mesuré ci-dessus a donc été obtenu \textbf{sans aucune
contrainte de préservation} : le compresseur pouvait jeter n'importe quoi sans que rien ne
proteste. Et \code{analyzer.rs:73} note que les sorties de l'analyseur « n'alimentent NI le router NI
le \code{must\_preserve} du compresseur » --- le seul composant capable de nommer les entités à
garder n'y est pas branché.
\end{alerte}

\begin{alerte}
\textbf{F2 --- La validation est circulaire.} \code{CompressorOutput::validate\_preservation}
compare \code{must\_preserve} à \code{self.preserved} --- la liste que le compresseur \textbf{déclare}
avoir préservée. Elle ne regarde \textbf{jamais} \code{self.summary}. Un compresseur qui rend un
résumé vide en déclarant \code{preserved: ["ACME", "2024-03-11"]} \textbf{passe la validation}.

C'est le cercle de \mm{M9} : on compare une affirmation à une exigence, jamais l'artefact à l'une ou
l'autre. \mm{M9} l'a cassé en lisant le jeton ; \mm{MZ} doit le casser en lisant le résumé.
\end{alerte}

\textbf{F3 --- Le modèle recrache parfois son gabarit.} \code{is\_degenerate\_summary} détecte une
sortie vide ou contenant \code{<the compressed text>} / \code{<source text>}. Ce détecteur n'existe
pas par précaution théorique : il existe parce que cela arrive. Son \emph{taux} n'a jamais été
chiffré, et c'est une des trois inconnues de \mm{MZ}.

\section{2 · Les trois quantités --- et non deux}

La feuille de route en demande deux : « le plafond du compresseur et le taux de préservation ». La
lecture du code en impose une troisième, et c'est celle qui décide.

\subsection{2.1 · Le plafond n'est pas un scalaire}

Comprimer davantage, c'est préserver moins. Le « plafond » n'est donc pas un nombre isolé mais le
point où la \textbf{courbe d'arbitrage} cesse d'être acceptable :

\begin{extrait}\small
\code{plafond := max\{ r : preservation\_verifiee(r) $\geq$ seuil \textbf{et} degenerescence(r) $\leq$ seuil \}}
\end{extrait}

où $r$ est le ratio effectif \code{compression\_ratio} (mots source / \code{word\_count}). Sans
seuil déclaré, « plafond » ne veut rien dire : §4 les fixe \emph{avant} l'expérience.

\subsection{2.2 · Deux taux de préservation, et leur écart}

\begin{center}\small
\renewcommand{\arraystretch}{1.25}
\begin{tabular}{@{}>{\raggedright\arraybackslash}p{3.4cm} >{\raggedright\arraybackslash}p{6.4cm} >{\raggedright\arraybackslash}p{6.0cm}@{}}
\toprule
& \textbf{Définition} & \textbf{Ce que sa mesure apprend} \\
\midrule
\textbf{Déclaré} & $|\,$\code{preserved} $\cap$ \code{must\_preserve}$\,| \;/\; |\,$\code{must\_preserve}$\,|$ & ce que le compresseur \emph{dit} de lui-même \\
\textbf{Vérifié} & entités retrouvées \textbf{dans \code{summary}} $/$ $|\,$\code{must\_preserve}$\,|$ & ce qui a \emph{réellement} survécu \\
\textbf{L'écart} & vérifié $-$ déclaré & si l'écart est non nul, \code{validate\_preservation} est \pire{décorative}, et toute garantie bâtie dessus est nulle \\
\bottomrule
\end{tabular}
\end{center}

L'écart est la quantité que personne n'a jamais regardée, et c'est la seule dont la valeur puisse
invalider un mécanisme plutôt qu'un réglage.

\subsection{2.3 · Pourquoi cela compte PLUS ici qu'ailleurs}

\begin{these}
\small Le Cadrage vit \textbf{à l'intérieur} d'un produit de preuve. La chaîne scelle fidèlement ce
qui a été enregistré --- donc si le compresseur laisse tomber une entité, le reçu prouve, avec toute
la rigueur cryptographique du dépôt, que \textbf{ce fait n'a jamais existé}. Le sceau serait parfait
et le contenu faux.

Ailleurs, perdre 5 \% des entités est une régression de qualité. Ici, c'est une
\textbf{falsification scellée}. C'est pourquoi les seuils de §4 sont plus durs qu'ils ne le seraient
pour un résumeur ordinaire.
\end{these}

\section{3 · Le protocole}

\subsection{3.1 · Ce qui rend MZ gratuit}

\code{models/compressor.gguf} est présent (911\,Mo). La mesure ne touche \textbf{aucune API
payante}, ne dépend d'aucun réseau, et n'exerce pas le pipeline --- seulement le compresseur. C'est
ce qui permet de la faire « avant d'en parler à quiconque », comme le dit la feuille de route.

\subsection{3.2 · Le corpus}

Un corpus \textbf{à entités connues}, construit et non trouvé : chaque document porte une liste
d'entités qu'on y a placées (identifiants de dossier, dates, montants, noms de fournisseurs,
clauses). Sans vérité de terrain, la préservation ne se mesure pas.

Trois familles, parce que le plafond n'a aucune raison d'être le même :
\begin{itemize}[leftmargin=1.2em,itemsep=1.5pt,topsep=2pt]
\item \textbf{dense} --- beaucoup d'entités par millier de mots (un dossier fournisseur) ;
\item \textbf{diffus} --- peu d'entités noyées dans de la prose (un compte rendu) ;
\item \textbf{adverse} --- entités en fin de document, ou répétées sous des formes voisines
      (« ACME S.A. », « ACME », « la société ») --- là où un résumeur extractif coupe en premier.
\end{itemize}

\subsection{3.3 · Le balayage}

Pour chaque document et chaque \code{target\_words} allant du généreux à l'agressif, on relève :

\begin{center}\small
\renewcommand{\arraystretch}{1.22}
\begin{tabular}{@{}>{\raggedright\arraybackslash}p{4.4cm} >{\raggedright\arraybackslash}p{11.4cm}@{}}
\toprule
\textbf{Mesure} & \textbf{Comment} \\
\midrule
ratio effectif & mots source $/$ \code{word\_count} --- \emph{effectif}, car \code{target\_words} est une consigne à $\pm$20 \%, pas un résultat \\
préservation déclarée & \code{preserved} $\cap$ \code{must\_preserve} \\
préservation \best{vérifiée} & recherche de chaque entité \textbf{dans le texte du résumé}, indépendamment de ce que le compresseur déclare \\
dégénérescence & \code{is\_degenerate\_summary} sur la sortie \\
déterminisme & deux passes identiques : la sortie est-elle stable ? \\
\bottomrule
\end{tabular}
\end{center}

\textbf{Un point de méthode.} La recherche « vérifiée » doit être \emph{stricte et déclarée} :
correspondance exacte, insensible à la casse, sur les frontières de mots. Toute tolérance
(lemmatisation, correspondance approchée) transformerait la mesure en jugement, et un jugement qu'on
règle soi-même donne toujours le résultat qu'on espère. Les formes voisines sont traitées par le
corpus \textbf{adverse}, pas par un assouplissement du contrôle.

\section{4 · Ce qui tue le Cadrage --- déclaré d'avance}

\begin{alerte}
Ces seuils sont écrits \textbf{avant} l'expérience. Les ajuster après avoir vu les chiffres
reviendrait à déplacer la cible après le tir --- exactement ce qu'on a refusé de faire au critère
« TLS » de \mm{M8}.
\end{alerte}

\begin{center}\small
\renewcommand{\arraystretch}{1.25}
\begin{tabular}{@{}>{\raggedright\arraybackslash}p{4.2cm} >{\raggedright\arraybackslash}p{11.6cm}@{}}
\toprule
\textbf{Verdict} & \textbf{Condition} \\
\midrule
\pire{TUÉ} & \textbf{L'écart déclaré/vérifié est non nul.} Le compresseur affirme préserver ce qu'il n'a pas préservé. Aucune valeur de ratio ne rachète cela : la garantie est un mensonge, et le pipeline la sceller\emph{ait}. \\
\pire{TUÉ} & Préservation \textbf{vérifiée} $< 99\,\%$ au ratio où l'économie devient intéressante. Dans un produit de preuve, une entité perdue sur cent est une falsification scellée, pas une imprécision. \\
\pire{TUÉ} & Dégénérescence $> 1\,\%$ à ce même ratio : le chemin de repli serait le chemin normal. \\
\midrule
\textbf{DÉCLASSÉ} & Le plafond existe mais l'économie qu'il autorise reste sous $\sim$20 \%. Avec $+79\,\%$ de latence déjà mesurés, le Cadrage ne paie pas --- il redevient une piste, pas un produit. \\
\midrule
\best{SAUVÉ} & Préservation vérifiée $\geq 99\,\%$ jusqu'à un ratio donnant $>$20 \% d'économie, écart nul, dégénérescence négligeable. \textbf{Le défaut est alors le CÂBLAGE}, pas le levier : \code{must\_preserve} vide (F1) devient un bogue à corriger, et non une limite à accepter. \\
\bottomrule
\end{tabular}
\end{center}

\begin{piste}
\small \textbf{L'issue la plus probable mérite d'être nommée.} Le compresseur peut très bien être
bon et le \emph{câblage} mauvais. C'est ce que suggère F1 : on n'a jamais demandé la préservation, on
ne peut donc pas savoir si elle tient. Dans ce cas \mm{MZ} ne tue rien --- il transforme une promesse
invérifiable en un défaut nommé, ce qui est déjà tout ce qu'on lui demande.
\end{piste}

\section{5 · Ce que MZ ne mesure pas}

\begin{itemize}[leftmargin=1.2em,itemsep=2pt,topsep=2pt]
\item \textbf{La qualité de la réponse finale.} \mm{MZ} mesure le compresseur, pas le pipeline. Un
      résumé qui préserve toutes les entités peut malgré tout mal servir le modèle frontière ; c'est
      une autre expérience, et elle coûte des API.
\item \textbf{La latence en charge.} Les $+79\,\%$ du rapport existant restent la référence ;
      \mm{MZ} relève le temps par appel, sans prétendre modéliser un service.
\item \textbf{Les autres robots.} Analyseur, vérifieur et reformulateur ont leurs propres plafonds.
      Les mesurer ensemble mélangerait les causes.
\item \textbf{La correction.} Comme partout dans cette série : ce qui est mesuré est ce qui a été
      \emph{préservé}, jamais ce qui est \emph{vrai}.
\end{itemize}

\section{6 · Livrables et ordre des gestes}

\begin{center}\small
\renewcommand{\arraystretch}{1.22}
\begin{tabular}{@{}>{\raggedright\arraybackslash}p{1.0cm} >{\raggedright\arraybackslash}p{7.4cm} >{\raggedright\arraybackslash}p{7.4cm}@{}}
\toprule
& \textbf{Geste} & \textbf{Pourquoi dans cet ordre} \\
\midrule
1 & Le corpus à entités connues, commité & sans vérité de terrain, rien n'est mesurable ; et un corpus écrit \emph{après} avoir vu les résultats n'est plus un corpus \\
2 & Le vérificateur \best{indépendant} de préservation & il ne doit rien devoir au compresseur : c'est lui qui casse le cercle de F2 \\
3 & Le balayage, et la courbe & la seule sortie qui compte \\
4 & Le verdict contre les seuils de §4 & écrits d'avance, appliqués sans discussion \\
5 & \code{scripts/mz-check.sh} & pour que la mesure se refasse, et qu'un chiffre cesse d'être une affirmation \\
\bottomrule
\end{tabular}
\end{center}

\textbf{Coût} --- 1,5 jour, aucune API, aucun réseau. Le seul coût réel est le temps de calcul local
du balayage.

\begin{these}
\small \textbf{Ce que MZ change, quel que soit son verdict.} Aujourd'hui le Cadrage est
\emph{exclu du produit} (feuille de route §4) sans que personne sache pourquoi au juste. Après
\mm{MZ}, il sera exclu \textbf{pour une raison chiffrée}, ou réintégré avec son plafond écrit. Les
deux valent mieux que l'état actuel --- et c'est pour cela que la feuille de route dit de le faire
\textbf{avant d'en parler à quiconque}.
\end{these}

\end{document}
